% ============================================================
% Secao 12: Conclusao
% ============================================================

\section{Conclus\~ao}
\label{sec:conclusao}

\subsection{Resumo das Contribui\c{c}\~oes}

Apresentamos o \aletheion{}, um LLM \textit{decoder-only} que integra
um sistema epist\^emico completo como m\'odulos neurais trein\'aveis.
As principais contribui\c{c}\~oes s\~ao:

\begin{enumerate}
    \item \textbf{Epistemologia intr\'inseca}: Pela primeira vez, um LLM
    produz tomografia epist\^emica completa (30+ campos) por token,
    treinada \textit{end-to-end} com o modelo de linguagem.

    \item \textbf{Manifold Riemanniano 5D}: Tokens s\~ao mapeados para
    um espa\c{c}o geom\'etrico aprendido onde dist\^ancias geod\'esicas
    definem confian\c{c}a calibrada.

    \item \textbf{Overhead m\'inimo}: Os 11 sub-heads epist\^emicos
    adicionam apenas ~364K par\^ametros (<0.3\%), demonstrando que
    capacidade epist\^emica n\~ao requer aumento significativo de
    par\^ametros.

    \item \textbf{Calibra\c{c}\~ao por constru\c{c}\~ao}: A loss MAD
    (\Cref{thm:calibration}) garante calibra\c{c}\~ao no equil\'ibrio,
    diferente de m\'etodos p\'os-hoc como temperature scaling.

    \item \textbf{Preven\c{c}\~ao de colapso}: O VI (\Cref{sec:vi})
    e EidosDecay (\Cref{sec:eidos}) previnem degenera\c{c}\~ao do
    manifold, um problema identificado empiricamente no sistema
    predecessor (ATIC).

    \item \textbf{Continual Learning nativo}: EWC + Replay integrados
    no trainer permitem treinamento em m\'ultiplas fases sem
    \textit{catastrophic forgetting}.

    \item \textbf{Escalabilidade}: Configura\c{c}\~oes de 1M a 640B
    par\^ametros com a mesma arquitetura, facilitando pesquisa
    (modelos pequenos) e deploy (modelos grandes).
\end{enumerate}

\subsection{Compara\c{c}\~ao com o ATIC}

O \aletheion{} surgiu da necessidade de superar as limita\c{c}\~oes do
sistema ATIC, um orquestrador determin\'istico que envolve LLMs externos:

\begin{table}[H]
\centering
\caption{Compara\c{c}\~ao ATIC vs \aletheion{}.}
\label{tab:vs_atic}
\begin{tabular}{lcc}
\toprule
\textbf{Aspecto} & \textbf{ATIC} & \textbf{\aletheion{}} \\
\midrule
Epist\^emico & Extr\'inseco & \textbf{Intr\'inseco} \\
Granularidade & Por resposta & \textbf{Por token} \\
Calibra\c{c}\~ao & Heur\'istica & \textbf{Via loss (BCE)} \\
Gradiente & Zero & \textbf{End-to-end} \\
Depend\^encia & LLM externo & \textbf{Auto-contido} \\
Lat\^encia & Pipeline multi-step & \textbf{1 forward pass} \\
Escalabilidade & CPU-bound & \textbf{GPU, 1M--640B} \\
Continual Learning & Nenhum & \textbf{EWC + Replay} \\
\midrule
Governan\c{c}a & Completa & Pendente \\
Agentes & 9+ registrados & Pendente \\
Flexibilidade & Troca LLM base & Fixo \\
\bottomrule
\end{tabular}
\end{table}

\subsection{Vis\~ao Futura: Integra\c{c}\~ao}

O cen\'ario ideal \'e a integra\c{c}\~ao do \aletheion{} como LLM base do ATIC,
combinando:

\begin{itemize}
    \item Epist\^emico intr\'inseco (\aletheion{}) com orquestra\c{c}\~ao (ATIC)
    \item Tomografia por token com governan\c{c}a e agentes
    \item Confian\c{c}a calibrada com \textit{policy binding}
    \item Manifold aprendido com navega\c{c}\~ao MPC em tempo real
\end{itemize}

Nesta configura\c{c}\~ao, o \aletheion{} de 354M funcionaria como um
avaliador epist\^emico das respostas de LLMs maiores, adicionando
tomografia calibrada a qualquer modelo externo.

\subsection{Trabalhos Futuros}

\begin{enumerate}
    \item \textbf{Avalia\c{c}\~ao emp\'irica}: Benchmarks de calibra\c{c}\~ao
    (ECE, MCE, Brier score) contra baselines (temperature scaling,
    ensemble, MC Dropout).

    \item \textbf{Transfer learning}: Verificar se o manifold 5D
    transfere entre dom\'inios (e.g., c\'odigo $\to$ texto cient\'ifico).

    \item \textbf{Escala}: Treinar vers\~oes 1.3B e 7B para validar
    que o sistema epist\^emico melhora com escala.

    \item \textbf{Integra\c{c}\~ao ATIC}: Substituir o LLM externo do
    ATIC pelo \aletheion{}, avaliando impacto na qualidade das respostas
    e na precis\~ao da tomografia.

    \item \textbf{Estudo de abl\c{c}\~ao}: Impacto individual de cada
    loss component e sub-head na calibra\c{c}\~ao final.

    \item \textbf{Efici\^encia do VI}: Comparar VI treinado
    \textit{end-to-end} vs VI determin\'istico (ATIC) em preven\c{c}\~ao
    de colapso.
\end{enumerate}

\subsection{C\'odigo e Reprodutibilidade}

O c\'odigo-fonte completo, incluindo todos os m\'odulos, testes (261),
configura\c{c}\~oes de escala e scripts de treinamento, est\'a dispon\'ivel
no reposit\'orio:

\begin{center}
    \texttt{github.com/gnai-creator/aletheion-llm-v2}
\end{center}

O modelo de 354M par\^ametros pode ser treinado em uma \'unica
RTX 4090 (24GB VRAM) em aproximadamente 8--10 dias com 7B tokens
do FineWeb-Edu.
